\documentclass[11pt, a4paper]{article}

% --- PACCHETTI ---
\usepackage[utf8]{inputenc}
\usepackage[italian]{babel}
\usepackage[margin=2.5cm]{geometry} % Imposta i margini della pagina
\usepackage{graphicx}               % Per inserire il logo
\usepackage{tabularx}               % Per tabelle larghe quanto la pagina
\usepackage[table]{xcolor}          % Per colorare le celle delle tabelle
\usepackage[hidelinks]{hyperref}    % Per stile link cliccabili
\usepackage{lmodern}      % Rende i font scalabili all'infinito
\usepackage{anyfontsize}  % Permette di usare dimensioni personalizzate
\usepackage{xurl}         % Permette ai link lunghi di non uscire dai margini
\urlstyle{same}           % Usa lo stesso font del documento

\begin{document}

% ==============================================================================
% FRONTESPIZIO (PAGINA 1)
% ==============================================================================
\begin{titlepage}
    \centering
    \vspace*{0.5cm}

    % LOGHI
    \includegraphics[width=0.45\textwidth]{assets/logo_unipd.jpg}
    \hfill
    \includegraphics[width=0.45\textwidth]{assets/logo.png}\par
    \vspace{0.5cm}

    % NOME DEL GRUPPO
    {\fontsize{45}{50}\sffamily\selectfont \textbf{On-By-One}\par} 
    \vspace{1cm}

    % TITOLO DEL DOCUMENTO
    {\huge \bfseries Analisi dei requisiti \par}
    
    \vspace{1.5cm}

    % MEMBRI DEL GRUPPO
    {\large
    \textbf{On-By-One}:\par
    \vspace{0.3cm}
    \begin{tabular}{l l}
        Sara Ballardin & 2101048 \\
        Shuhui Chen & 2138001 \\
        Simone Giuliano & 2145591 \\
        Alvise Greggio & 2141004 \\
        Daniele Nuzzaci & 2139039 \\
        Federico Temporin & 2146764 \\
        Tommaso Zorzetto & 2148629 \\
    \end{tabular}
    \par}

    \vfill

    % TABELLA INFORMAZIONI DOCUMENTO E DESTINATARI
    \begin{table}[h]
        \centering
        \arrayrulecolor{gray!50}
        \renewcommand{\arraystretch}{1.5}
        \begin{tabularx}{\textwidth}{|c|c|c|X|}
            \hline
            \rowcolor{gray!10} \multicolumn{4}{|c|}{\textbf{Informazioni documento}} \\
            \hline
            \rowcolor{gray!10} \textbf{Versione} & \textbf{Data} & \textbf{Stato} & \multicolumn{1}{c|}{\textbf{Destinatari}} \\
            \hline
            - & - & Bozza & 
            On-By-One \newline 
            Prof. Tullio Vardanega \newline 
            Prof. Riccardo Cardin \newline 
            Vimar S.P.A. \\
            \hline
        \end{tabularx}
        \arrayrulecolor{black}
    \end{table}
    
    \vspace{0.5cm}
    
    % CONTATTI
    {\large Contatti: \href{mailto:onbyone.swe@gmail.com}{onbyone.swe@gmail.com} }
    
\end{titlepage}

\newpage


% ==============================================================================
% REGISTRO DELLE MODIFICHE (PAGINA 2)
% ==============================================================================
\section*{Registro delle modifiche}

\vspace{1cm}
\begin{center}
    \arrayrulecolor{gray!50}
    \renewcommand{\arraystretch}{1.5}
    \begin{tabularx}{\textwidth}{|c|c|c|X|c|}
        \hline
        \rowcolor{gray!10} 
        \multicolumn{1}{|c|}{\textbf{Versione}} & 
        \multicolumn{1}{c|}{\textbf{Data}} & 
        \multicolumn{1}{c|}{\textbf{Autore}} & 
        \multicolumn{1}{c|}{\textbf{Descrizione}} & 
        \multicolumn{1}{c|}{\textbf{Verificatore}} \\
        \hline
        - & - & - & - & - \\
        \hline
        0.1.2 & 2026/11/14 & Shuhui Chen & Stesura paragrafo Scopo e obiettivi del prodotto & - \\
        \hline
        0.1.2 & 2026/11/14 & Simone Giuliano & Stesura paragrafi 3.1 e 3.2 & - \\
        \hline
        0.1.1 & 2026/11/13 & Sara Ballardin & Stesura dell'introduzione & - \\
        \hline
        0.1.0 & 2026/11/09 & Sara Ballardin & Definizione struttura del documento & - \\
        \hline
    \end{tabularx}
    \arrayrulecolor{black}
\end{center}
\vfill
\newpage


% ==============================================================================
% INDICE (PAGINA 3)
% ==============================================================================
\tableofcontents
\newpage


% ==============================================================================
% 1. INTRODUZIONE
% ==============================================================================
\section{Introduzione}

\subsection{Scopo del documento}
Il documento di \textbf{Analisi dei requisiti} raccoglie e descrive le funzionalità e i vincoli del progetto \textbf{AURA} (AI for Understanding, Reasoning and Automation), proposto da Vimar S.P.A. e sviluppato dal gruppo On-By-One.

\noindent L'obiettivo è trasformare le richieste del capitolato e quanto emerso negli incontri con l'azienda in un elenco chiaro e verificabile di requisiti software. Partendo dallo studio degli attori e dei casi d'uso in UML, il documento definisce cosa deve fare il sistema per soddisfare le esigenze dell'utente, lasciando le scelte su come realizzarlo alla successiva fase di progettazione.

\noindent Il documento rappresenta il punto di riferimento per tutti i soggetti coinvolti nel progetto:
\begin{itemize}
    \item \textbf{Proponente (Vimar S.P.A.) e committenti (Prof.\ Tullio Vardanega e Prof.\ Riccardo Cardin):} per verificare che le richieste iniziali siano state comprese correttamente e nella loro interezza;
    \item \textbf{Progettisti e programmatori:} come guida per definire l'architettura del sistema e sviluppare le funzionalità previste;
    \item \textbf{Verificatori:} come base per scrivere i test e controllare che il prodotto realizzato rispetti tutti i requisiti fissati.
\end{itemize}

\subsection{Glossario}
Per evitare ambiguità e facilitare la lettura di termini tecnici, acronimi e parole con un significato specifico nel progetto, il gruppo ha redatto un documento interno dedicato chiamato \textbf{\href{https://on-by-one-team.github.io/On-By-One.github.io/docs/pdf/glossario/glossario.pdf}{Glossario}}, che raccoglie tutte le relative definizioni.

\subsection{Riferimenti}
Di seguito sono elencati i documenti vincolanti di progetto, le regole interne del gruppo e le fonti informative (standard internazionali e materiale del corso) consultate per redigere questo documento.

\subsubsection{Riferimenti normativi}
\begin{itemize}
    \item Capitolato d'appalto AURA (Vimar S.P.A.) \\
    \url{https://www.math.unipd.it/~tullio/IS-1/2026/Progetto/C1-AURA.pdf} \\
    Ultimo accesso: 13 novembre 2026.

    \item Regolamento del progetto didattico \\
    \url{https://www.math.unipd.it/~tullio/IS-1/2026/Dispense/PD1.pdf} \\
    Ultimo accesso: 13 novembre 2026.

    \item Norme di progetto (documento interno del gruppo On-By-One) \\
    \url{https://on-by-one-team.github.io/On-By-One.github.io/docs/pdf/norme-di-progetto-documentazione.pdf} \\
    Ultimo accesso: 13 novembre 2026.
\end{itemize}

\subsubsection{Riferimenti informativi}
\begin{itemize}
    \item Analisi dei requisiti (Prof.\ Tullio Vardanega) \\
    \url{https://www.math.unipd.it/~tullio/IS-1/2026/Dispense/T05.pdf} \\
    Ultimo accesso: 13 novembre 2026.

    \item Diagrammi dei casi d'uso (Prof.\ Riccardo Cardin) \\
    \url{https://www.math.unipd.it/~rcardin/swea/2022/Diagrammi%20Use%20Case.pdf} \\
    Ultimo accesso: 13 novembre 2026.
\end{itemize}

\newpage


% ==============================================================================
% 2. DESCRIZIONE DEL PRODOTTO
% ==============================================================================
\section{Descrizione del prodotto}

\subsection{Scopo e obiettivi del prodotto}
Il progetto consiste nella realizzazione di AURA (\emph{AI for Understanding, Reasoning and Automation}), progettata per l’integrazione avanzata tra l'impianto domotico \emph{View Wireless} e i principali servizi di intelligenza artificiale generativa. 
 Il suo scopo è mostrare la fattibilità e le potenzialità
dell'integrazione tra l'impianto domotico \emph{View Wireless} e i principali
servizi di intelligenza artificiale generativa, sfruttando lo standard aperto \emph{Model Context Protocol} (MCP) e le \emph{KNX IoT 3rd-party API} (v2.1.0).
L'idea di fondo è usare servizi GenAI già esistenti come interfaccia in linguaggio naturale verso la smart home, senza sviluppare un assistente AI dedicato. Il prodotto non intende sostituire le applicazioni esistenti (come l'app Vimar View), ma affiancarle come canale complementare.

Il prodotto deve soddisfare i seguenti obiettivi:
\begin{itemize}
    \item realizzare un'infrastruttura \textit{}{Cloud}-ready \textbf{containerizzata}, definita come \textit{Infrastructure as Code} e compatibile con AWS;
    \item realizzare un \textbf{server MCP} e un \textbf{client MCP} che permettano \textbf{lettura} e \textbf{scrittura} sui dispositivi, ricostruendo le relazioni tra\textit{Location}, \textit{Function} e \textit{Datapoint};
    \item garantire un accesso \textbf{sicuro} all'impianto tramite \textit{OAuth2} e account MyVimar;
    \item collegare il server a un \textbf{servizio AI commerciale} noto (Claude, ChatGPT, Copilot, Gemini);
    \item consegnare software \textbf{open source}, documentato e testato, con copertura minima del 75\% per unità e integrazione e dell'80\% per i test end-to-end.
\end{itemize}

\subsection{Contesto d'uso e prospettiva del prodotto}
AURA si colloca come strato di integrazione tra il servizio AI scelto dall'utente e il \textbf{Cloud Vimar}. Il server MCP, ospitato nell'infrastruttura Cloud-ready, fa da \textbf{adattatore} tra le due parti: riceve le richieste del servizio AI tramite \textit{client MCP}, le traduce in chiamate alle KNX IoT 3rd-party API, autenticate con OAuth2, e restituisce i risultati. A sua volta il Cloud Vimar comanda i dispositivi tramite \textbf{Gateway BLE} e rete \textit{Bluetooth Mesh}. Il flusso completo, dalla richiesta dell'utente alla risposta, è quindi il seguente:
\begin{enumerate}
    \item l'utente invia una richiesta al \textbf{servizio AI esterno};
    \item il servizio AI interroga il \textbf{server MCP} tramite client MCP;
    \item il server MCP chiama le \textbf{KNX IoT 3rd-party API};
    \item il Cloud Vimar esegue il comando sui dispositivi o ne legge lo stato, e il risultato risale lungo la stessa catena.
\end{enumerate}

\subsection{Funzionalità principali}
Le funzionalità \textbf{obbligatorie} sono:
\begin{itemize}
    \item \textbf{autenticazione} con account MyVimar tramite OAuth2 e
    \textbf{selezione dell'impianto}, mantenuta per tutta la conversazione;
    \item \textbf{esplorazione} di ambienti, funzioni e dispositivi;
    \item \textbf{lettura dello stato} per dispositivo, ambiente, tipologia di funzione o intero impianto;
    \item \textbf{scrittura} su dispositivi singoli, multipli o sull'intero impianto;
    \item \textbf{gestione del linguaggio naturale}, comprese richieste ambigue, sinonimi, richieste implicite e conferma per i comandi estesi;
    \item \textbf{gestione degli errori}, con indicazione delle alternative disponibili;
    \item \textbf{avvio con un solo comando} (\textit{docker compose up}) e configurazione tramite file di variabili d'ambiente.
\end{itemize}

A queste si aggiungono funzionalità \textbf{opzionali}, che estendono il prodotto senza essere indispensabili: il supporto all'endpoint \emph{Device} delle KNX IoT 3rd-party API, l'uso del client MCP come estensione VSCode o da CLI/TUI, il rilascio effettivo su AWS e una copertura dei test superiore al 90\%. Infine, le funzionalità \textbf{nice-to-have} riguardano richieste su più \emph{impianti}, storico e operazioni sui dati, azioni condizionate, scenari composti, trigger ed eventi 

\subsection{Caratteristiche degli utenti}
L'\textbf{utente residenziale} è il principale destinatario del prodotto. Dispone di un \emph{impianto} View Wireless e di un account MyVimar, e comanda la propria abitazione in linguaggio naturale tramite il servizio AI che preferisce.
\newpage


% ==============================================================================
% 3. CASI D'USO (USE CASES)
% ==============================================================================
\section{Casi d'uso}

\subsection{Introduzione ai casi d'uso}
Un \textit{caso d'uso} rappresenta una descrizione chiara e dettagliata di un insieme di \textit{scenari} che condividono il medesimo obiettivo per un \textit{attore}. Tali \textit{scenari} descrivono il modo in cui un \textit{attore} interagisce con il sistema e le funzionalità che quest'ultimo gli fornisce.

 Al fine di garantire l'identificazione e la tracciabilità delle funzionalità, la codifica adottata è \texttt{UC[Codice\_Padre].[Codice\_Figlio]}. I \textit{sotto-casi d'uso} ereditano l'identificativo del \textit{caso d'uso} principale da cui derivano. Ogni livello gerarchico successivo è separato dal precedente tramite il carattere punto (\texttt{.}), definendo così una struttura ad albero.
Ad integrazione della descrizione testuale, ogni \textit{caso d'uso} è accompagnato da un diagramma UML 2.0 che illustra le interazioni previste tra gli \textit{attori} coinvolti e le funzionalità del sistema.

\subsection{Formato di specifica}
All'interno del documento, ogni \textit{caso d'uso} verrà descritto attraverso le seguenti caratteristiche, da intendersi secondo le definizioni di seguito riportate:

\begin{itemize}
    \item \textbf{Attori principali:} Rappresentano i ruoli assunti dalle entità che interagiscono in modo attivo con il sistema, al fine di raggiungere il proprio obiettivo.
    
    \item \textbf{Attori secondari:} Rappresentano eventuali entità esterne coinvolte passivamente nell'interazione per supportare o completare l'esecuzione del \textit{caso d'uso}.
    
    \item \textbf{Pre-condizioni:} L'insieme delle condizioni di sistema che devono risultare verificate affinché il \textit{caso d'uso} possa iniziare.
    
    \item \textbf{Post-condizioni:} Le condizioni finali che il sistema garantisce al termine dell'esecuzione del \textit{caso d'uso}. Per ogni \textit{scenario} (principale o alternativo) vengono definite specifiche post-condizioni; al termine dell'interazione risulteranno verificate quelle relative allo \textit{scenario} effettivamente percorso.
    
    \item \textbf{Scenario principale:} Univoca sequenza ordinata di azioni che rappresenta il percorso standard, privo di errori o eccezioni, che conclude con il completamento del \textit{caso d'uso}.
    
    \item \textbf{Scenari alternativi:} Sequenze di eventi (da zero a molteplici) che modellano eccezioni, errori o deviazioni dal flusso standard.
    
    \item \textbf{Inclusioni:} Riferimenti ad altri \textit{casi d'uso} che devono obbligatoriamente essere eseguiti per poter concludere il \textit{caso d'uso} includente.
    
    \item \textbf{Estensioni:} Riferimenti a \textit{casi d'uso} estendenti che vengono eseguiti solo al verificarsi di specifiche condizioni, che aumentano o alterano le funzionalità in uno \textit{scenario} alternativo.

    \item \textbf{Generalizzazioni:} Relazioni in cui un \textit{caso d'uso} derivato eredita interamente le caratteristiche di un \textit{caso d'uso} base, potendone specializzare il comportamento.
\end{itemize}
 
\subsection{Classificazione degli attori}
\subsubsection{Attori principali}
Elencare e descrivere gli attori che avviano proattivamente le interazioni con il sistema per raggiungere un obiettivo di business, allegando il diagramma generale delle relazioni tra attori.
 
\subsubsection{Attori secondari}
Elencare e descrivere i sistemi, le API o i servizi esterni che prendono parte all'esecuzione dei casi d'uso senza avviarli direttamente.

\subsection{Elenco dei casi d'uso}
Raccogliere i casi d'uso raggruppati per attore principale o per area funzionale. Per ognuno inserire diagramma UML e descrizione secondo il formato definito sopra.

\newpage


% ==============================================================================
% 4. REQUISITI
% ==============================================================================
\section{Requisiti}

\subsection{Classificazione e codifica dei requisiti}
Descrivere la convenzione di codifica dei requisiti (tipologia, numero progressivo, priorità obbligatorio/desiderabile/opzionale) e le fonti da cui derivano (casi d'uso, capitolato, verbali degli incontri con la proponente).

\subsection{Requisiti funzionali}
Requisiti che descrivono cosa il sistema deve fare (codice identificativo, descrizione e fonte (caso d'uso)).

\subsection{Requisiti di qualità}
Requisiti sulla qualità di prodotto e di processo (testing, copertura, documentazione da consegnare, manutenibilità).

\subsection{Requisiti di vincolo}
Requisiti che impongono limiti alle scelte progettuali e tecnologiche.

\subsection{Requisiti prestazionali}
Requisiti misurabili su tempi di risposta, capacità e scalabilità, se presenti.

\newpage


% ==============================================================================
% 5. TRACCIAMENTO DEI REQUISITI
% ==============================================================================
\section{Tracciamento}

\subsection{Tracciamento fonti -- requisiti}
Tabella che mappa ciascuna fonte formale (capitolato d'appalto, verbali di riunione esterni, casi d'uso) ai requisiti che ne sono scaturiti.

\subsection{Tracciamento casi d'uso -- requisiti}
Tabella di correlazione che associa ogni singolo caso d'uso ai relativi requisiti funzionali soddisfatti.


% ==============================================================================
% 6. RIEPILOGO
% ==============================================================================
\section{Riepilogo dei requisiti}
Tabella riassuntiva con il conteggio complessivo dei requisiti per tipologia (funzionali, qualitativi, di vincolo, prestazionali) e per grado di priorità (obbligatori, desiderabili, opzionali).

\end{document}
